\chapter{Kết luận và Hướng phát triển}
\label{ch:conclusion}

Chương này tổng kết các kết quả đã đạt được của khóa luận, nêu ý nghĩa và hạn chế, đồng thời đề
xuất các hướng phát triển trong tương lai.

\section{Tổng kết}
\label{sec:summary}

Khóa luận đã nghiên cứu và xây dựng một hệ thống trả lời câu hỏi có ràng buộc về thời gian cho
\textbf{tiếng Việt}, dựa trên phương pháp Suy luận Trừu tượng Quy nạp (ARI). Xuất phát từ hai hạn
chế cốt lõi của mô hình ngôn ngữ lớn -- thiếu tri thức thời gian và thiếu năng lực suy luận thời
gian phức tạp -- khóa luận kế thừa kiến trúc hai pha của ARI, tách bạch khâu tương tác tri thức khỏi
khâu suy luận trừu tượng, qua đó cung cấp chỗ dựa tri thức xác thực cho mô hình đồng thời giảm thiểu
nhiễu.

Cụ thể, khóa luận đã hoàn thành ba đóng góp chính:

\begin{enumerate}
  \item \textbf{Bộ dữ liệu CronQ-VN.} Khóa luận thiết kế một quy trình hoàn chỉnh trích xuất đồ thị
  tri thức thời gian tiếng Việt từ Wikidata (khoảng 238 nghìn sự kiện, 37,5 nghìn thực thể, 271
  quan hệ) và sinh tự động các câu hỏi kèm đáp án chuẩn, tạo nên bộ dữ liệu TKGQA tiếng Việt đầu
  tiên với năm loại câu hỏi thời gian.
  \item \textbf{Hệ thống ARI tiếng Việt vận hành cục bộ.} Khóa luận cài đặt lại toàn bộ khung ARI
  cho tiếng Việt, vận hành pha suy luận hoàn toàn bằng mô hình ngôn ngữ lớn chạy cục bộ thay cho API
  thương mại, kèm cơ chế liên kết thực thể tiếng Việt cho câu hỏi mở.
  \item \textbf{Bộ lọc hành động dựa trên cross-encoder.} Khóa luận chỉ ra hạn chế của bộ lọc hành
  động bằng độ tương đồng cosine trong ARI gốc và đề xuất thay thế bằng một bộ lọc nhận biết ngữ
  cảnh dựa trên cross-encoder, mã hóa đồng thời hành động và lịch sử suy luận.
\end{enumerate}

Kết quả thực nghiệm trên CronQ-VN \todo{tóm tắt kết quả chính sau khi chạy đánh giá: độ chính xác
tổng của ARI, mức cải thiện so với LLM trực tiếp và các baseline, cùng kết luận về vai trò của hướng
dẫn trừu tượng và bộ lọc cross-encoder.}

\section{Ý nghĩa}
\label{sec:significance}

Về mặt học thuật, khóa luận đóng góp \emph{bộ dữ liệu và hệ thống TKGQA tiếng Việt đầu tiên}, tạo
tiền đề cho các nghiên cứu tiếp theo về hỏi đáp thời gian trong tiếng Việt -- một ngôn ngữ còn ít
tài nguyên cho bài toán này. Việc vận hành toàn bộ pha suy luận bằng mô hình ngôn ngữ lớn cục bộ cho
thấy tính khả thi của các hệ thống hỏi đáp tri thức \emph{không phụ thuộc dịch vụ thương mại}, phù
hợp với các bối cảnh đòi hỏi chi phí thấp hoặc bảo mật dữ liệu.

Về mặt phương pháp, khóa luận củng cố một hướng tiếp cận giàu tiềm năng: kết hợp \emph{bộ nhớ kinh
nghiệm} với mô hình ngôn ngữ lớn, để mô hình chủ động kiến tạo và tái sử dụng các chiến lược suy
luận trừu tượng thay vì chỉ tiếp nhận thông tin thụ động. Về mặt ứng dụng, hệ thống có thể phục vụ
các tác vụ thực tế như tra cứu sự kiện lịch sử, dựng tiểu sử nhân vật và phân tích dữ liệu theo dòng
thời gian.

\section{Hạn chế}
\label{sec:limitations}

Bên cạnh các kết quả đạt được, khóa luận còn một số hạn chế. Thứ nhất, chất lượng của phương pháp
luận trừu tượng phụ thuộc vào năng lực của mô hình ngôn ngữ lớn; khi chạy cục bộ với mô hình quy mô
vừa phải, hướng dẫn sinh ra có thể chưa thật sắc bén. Thứ hai, do dựa trên suy luận nhiều bước, hệ
thống có chi phí và thời gian suy luận cao hơn so với trả lời trực tiếp. Thứ ba, phạm vi hiện tại
giới hạn ở suy luận thời gian, với bộ dữ liệu CronQ-VN ở hạt thời gian cấp năm và một tập quan hệ
tuyển chọn, nên chưa bao quát hết độ phức tạp của ngôn ngữ tự nhiên. Cuối cùng, đóng góp về bộ lọc
cross-encoder mới ở mức đề xuất và đang trong quá trình hoàn thiện thực nghiệm.

\section{Hướng phát triển}
\label{sec:future}

Từ các hạn chế trên, khóa luận đề xuất một số hướng phát triển.

\paragraph{Hoàn thiện và đánh giá bộ lọc cross-encoder.} Hướng trước mắt là hiện thực hóa đầy đủ bộ
chấm điểm cross-encoder, xây dựng tập dữ liệu huấn luyện từ các trace lịch sử với chiến lược lấy mẫu
âm phù hợp, rồi đánh giá định lượng mức cải thiện so với bộ lọc cosine, qua đó kiểm chứng giả thuyết
ở Mục~\ref{sec:ce-analysis}.

\paragraph{Nâng cao liên kết thực thể tiếng Việt.} Có thể tích hợp hoàn chỉnh quy trình nhận dạng
thực thể bằng mô hình NER tiếng Việt kết hợp tìm kiếm tương đồng trong cơ sở dữ liệu vector, nhằm
xử lý tốt hơn các câu hỏi mở và giảm nhóm lỗi liên kết thực thể đã chỉ ra ở Mục~\ref{sec:errors}.

\paragraph{Mở rộng bộ dữ liệu.} Bộ dữ liệu CronQ-VN có thể được mở rộng theo nhiều chiều: bổ sung
thêm quan hệ và lĩnh vực, hỗ trợ nhiều \emph{hạt thời gian} hơn (tháng, ngày), tăng quy mô câu hỏi,
và đa dạng hóa các dạng câu hỏi thời gian phức tạp hơn nữa.

\paragraph{Tăng cường năng lực suy luận.} Các hướng khả dĩ gồm: tinh chỉnh một mô hình ngôn ngữ lớn
tiếng Việt riêng cho bài toán nhằm cải thiện chất lượng quyết định và độ tuân thủ định dạng; khảo
sát các chuỗi suy luận sâu hơn; và cân bằng tốt hơn giữa độ sâu suy luận và chi phí tính toán.

\paragraph{Mở rộng phạm vi bài toán.} Về lâu dài, phương pháp có thể được mở rộng ra ngoài miền suy
luận thời gian thuần túy, chẳng hạn suy luận không -- thời gian (\textit{spatio-temporal}) hoặc các
miền suy luận khác, tận dụng tính tách rời thấp giữa hệ thống và tập thao tác nguyên tử để dễ dàng
bổ sung thao tác mới cho từng miền.

\bigskip

Tóm lại, khóa luận đã xây dựng nền tảng đầu tiên cho bài toán hỏi đáp có ràng buộc thời gian trên đồ
thị tri thức tiếng Việt, bao gồm cả bộ dữ liệu, hệ thống và một đề xuất cải tiến. Những kết quả và
hướng phát triển nêu trên hy vọng sẽ là cơ sở hữu ích cho các nghiên cứu tiếp theo trong lĩnh vực
này.
